\chapter{Reinforcement Learning and Graph Neural Networks (GNNs)}
In this chapter we introduce the framework on which the quasioptimal solution
of finding the arbitrage walk is going to be based on. We are going to give a
brief introduction to Monte Carlo based Reinforcement Learning which we will
apply based on a model of the environment of a network of CFMM. This chapter
is largely based on \cite[Chapter17]{ml_foundations} for the foundations of
Reinforcement Learning and \cite{sutton_barto} for the REINFORCE based MCTS
algorithm.
\section{Markov Decision Process}
Generally in comparison to supervised learning, reinforcement learning is not
trained directly with labeled data, instead it learns through an agent which takes
actions and is rewarded based on the outcome of these actions. The idea is
to, in expectation, increase the rewards of the agent throughout the learning
process by `trial and error'. Hence, reinforcement learning
is better suited for problems with large solution spaces and decision making
setting, where we have dynamical environments. In these cases the
combinatorial complexity of these problems is too big, such that supervised
learning would lack labeled data. For instance the Shannon number is
$10^{120}$, which is the lower bound of all possible positions in chess.
In this chapter we will go through the foundation of reinforcement learning
based on \cite[Chapter~17]{ml_foundations}.

The general the learning scenario is based on the agent collecting
information through actions, interacting with the environment. In response to
an action the agent will enter a new state and will be given a reward
depending on the general goal \ref{fig: 1}. For example the agent is rewarded
a score value of $+1$ if set of actions leads to a win in a decision based
game. At each state the agents objective is to maximize the expected reward,
which translates to taking in expectation the best action at a given state,
called the policy.
\begin{figure}[H]
    \centering
    \begin{tikzpicture}

        \node[align=center, ellipse, draw, ultra thick] (a)  {Agent};
        \node[align=center, ellipse, draw, below=3cm of a, ultra thick] (e) {Environment};


        \draw[->, ultra thick, blue] (a) to[out=0, in=0] node [midway,
            right, text=black] {action} (e);
        \draw[->, ultra thick, red] (e) to[out=160, in=200]
            node [midway, right, text=black] {state} (a) ;
        \draw[->, ultra thick, red] (e) to[out=180, in=180]
            node [midway, left, text=black] {reward} (a);

    \end{tikzpicture}
    \caption{Reinforcement learning scenario \label{fig: 1}}
\end{figure}

At each step the agent faces a decision to choose an action leading to a new
state and collect reward associated with it. Yet the agent is faced with a
dilemma of whether to maybe choose an action to an unexplored state which
would potentially lead to a higher reward or stick with an already
known-to-work set of actions. In reinforcement learning this is called
`exploration versus exploitation' trade off and is manifested and controlled
in the scope of the learning algorithm.

Also, a key difference to traditional supervised learning, reinforcement
learning does not sample features based on a fixed distribution. Essentially
the training and testing phases are not separated, but rather combined in
each step. The two main learning settings are distinguished, the setting
where the full environmental model is known and the one where the
environmental model is unknown. The mathematical model which describes the
agent's iteration with the environment is called the Markov Decision Process (MPD).
\begin{definition}[Markov decision process]
    The MDP is defined by $4$ tuple $(S, A, \delta, R_a)$:
    \begin{itemize}[nosep]
        \item The set of all states $S$, with a start state $s_0 \in S$,
        \item The set of all actions $A$,
        \item The transition probability $\mathbb{P}[s'|s, a] = \int_{S'}
            P_a(s, s') ds'$
        \item The reward probability $\mathbb{P}[r'|s, a] =
            \int_{S'}R_a\left(s', s  \right)ds' $
    \end{itemize}
\end{definition}
The transition probability $\mathbb{P}[s'|s, a]$ describes the probability of
the agent landing in state $s'$, when being initially at state $s$ and takes
an action $a$. Similarity the reward probability is explained in the same
way. The structure of state transition only depends on the last state and the
action taken and not on the whole history of states and actions hence the
Markov property and the naming of the model. Usually and in our case, the
state transitions and rewards are deterministic, in this case the transition
and reward probability is $1$ and the MDP is called deterministic with the
reward being $r(s, a)$.
\newpage
\begin{figure}[H]
    \centering
    \begin{tikzpicture}

        \node[align=center, circle, draw, ultra thick, minimum size=1cm] (1)  {$s_t$};
        \node[align=center, circle, draw, ultra thick, right=of 1, minimum size=1cm] (2)
            {$s_{t+1}$};
        \node[align=center, circle, draw, ultra thick, right=of 2, minimum
            size=1cm] (3)
            {$s_{t+2}$};

        \draw[->, dotted, ultra thick, -latex] (-1.5, 0) to (1);
        \draw[->, dotted, ultra thick, -latex] (3) to (5.5, 0);
        \draw[->, -latex] (1) to node[midway, below] {$a_t$} node[midway,
            above] {$r_{t+1}$} (2);
        \draw[->, -latex] (2) to node[midway, below] {$a_{t+1}$} node[midway,
            above] {$r_{t+2}$} (3);

    \end{tikzpicture}
    \caption{Deterministic discrete MDP \label{fig: 2} \cite{ml_foundations}}
    \label{}
\end{figure}
As per the discrete model, the agent preforms actions throughout a set of
decision epochs $t \in {0,\ldots,T}$ \ref{fig: 2}. This is the usual model
adopted for a wide range of problems. The MDP is called finite when both $S$
and $A$ are finite sets. Furthermore, the MDP has a finite time horizon when
$T$ is finite.
\newline
Generally we would like to find a to find a good policy for the decisions,
i.e. probabilities at each state telling us which action to take. The policy
is a function of states and returns a probability distribution over
the possible actions from the given state. We choose the action with the
higher probability in the probability distribution given by the policy.
\begin{definition}[Policy]
    A policy is a mapping $\pi : S \to \Delta\left( A \right)$, where $S$ is
    the state space and $\Delta(A)$ is the set of probability distributions
    over A, the action space.

    We call the policy deterministic if:
    \begin{align}
        \forall s \in S\;\; \exists !  a\in A:\quad \pi(a|s) = 1,
    \end{align}
    where $\pi(a|s)$ is the probability of landing in action $a$ given state
    $s$, the notation `$|$' serves as a notation for a probability
    distribution over $a \in A(s)$ for all $s \in S$. In deterministic case
    we have that $\text{im}(\pi) = A$, and we can say $\pi: S \to A$ with
    $\pi(s) = a$.
\end{definition}
Together with the definition of a policy, the aim of the agent can be
expressed as trying to find a policy that maximizes the total in
return in expectation. For a discrete policy the total return is
\begin{align}
    &\sum_{t=0}^{T} r(s_t, \pi(s_t)) & \text{finite horizon}\\
    &\sum_{t=0}^{\infty}\gamma^{t} r(s_t, \pi(s_t)) & \text{infinite horizon}
\end{align}
where $r(s_t, \pi(s_t))$ is the reward function for the next state and
$\gamma \in [0, 1)$ a discount factor for future rewards. To quantify what
the optimal policy is, we define the value of the policy at a given state.
\begin{definition}[Policy value]
    The value $V_\pi(s)$ of the policy $\pi$ at state $s \in S$ is defined as
    the expected value of the total reward over the policy distribution
    $\pi(s_t)$ of state $s_t \in S$ at epoch $t \le T$. For the finite horizon
    it is
    \begin{align}
        V_\pi(s) = \mathbb{E}_{a_t\sim\pi(s_t)}\left[\sum_{t=0}^{T}r(s_t,
        a_t) \Bigg| s_0 = s \right].
    \end{align}
    For the infinite horizon for $\gamma \in [0, 1)$
    \begin{align}
        V_\pi(s) =
        \mathbb{E}_{a_t\sim\pi(s_t)}\left[\sum_{t=0}^{\infty}\gamma^{t} r(s_t,
        a_t) \Bigg| s_0 = s \right].
    \end{align}
\end{definition}
Thereby the optimal policy can be defined in a standard way
\begin{definition}[Optimal policy]
    The policy $\pi^{*}$ is called optimal if for any policy $\pi$ it
    holds that
    \begin{align}
        V_{\pi^{*}}(s) \ge V_\pi(s) \quad \forall s \in S.
    \end{align}
\end{definition}
Similarly to the policy value we can quantify the pair $(s, a)$ associated to a
policy $\pi$ by a value, the so called state-action value
\begin{definition}[State-Action value]
    \label{def: state-action}
    The state-action value function $Q_\pi$ associated with a policy $\pi$ is
    defined on all state-action pairs $(s, a) \in S \times A$
    \begin{align}
        Q_\pi(s, a) &= \mathbb{E}[r(s,a)]
        \mathbb{E}_{a_t\sim\pi(s_t)}\left[ \sum_{t=1}^{\infty}\gamma^{t}
            r(s_t, a_t)\Bigg |\; s_0 = s,\; a_0 = a \right] \\
        \nonumber
            &= \mathbb{E}\left[r(s,a) + \gamma V_\pi(s_1)\Big| \;s_0=s,\;
            a_0=a   \right].
    \end{align}
\end{definition}
Note that the expected value of the state-value function at $(s, a)$ over the policy
distribution at state $s$ is the policy value at $s$, i.e.
\begin{align}
    \label{eq: value_expectation}
    \mathbb{E}_{a\sim\pi(s)}\left[Q_\pi(s, a)\right] = V_\pi(s).
\end{align}
Now that the framework for is set we can show that there exists an optimal
policy
\begin{theorem}[Policy Improvement Theorem]
    For any two policies $\pi'$ and $\pi$, it holds that
    \begin{align}
         &\mathbb{E}_{a\sim\pi'(s)}\left[Q_\pi(s, a)\right]
        \ge \mathbb{E}_{a\sim\pi(s)}\left[Q_\pi(s, a)\right]  \quad \forall s \in S
        \\
         &\quad \implies  V_{\pi'}(s) \ge V_\pi (s) \quad
         \forall s \in S\nonumber.
    \end{align}
\end{theorem}
\begin{theorem}[Bellman's optimality condition]
    A policy $\pi$ is optimal if and only if for all state-action pairs
    $(s, a) \in S x A$ with $\pi(s)(a) > 0$ it holds that
    \begin{align}
        a \in \argmax_{a'\in A} Q_\pi(s, a').
    \end{align}
\end{theorem}
\begin{theorem}{Existence of optimal deterministic policy}
    Any finite MDP has an optimal deterministic policy.
\end{theorem}
Now we consider an optimal deterministic policy $\pi^{*}$ with its policy value and
state-action value functions $V^{*}$ and $Q^{*}$. The optimal policy for all
states is
\begin{align}
    \pi^{*}(s) = \argmax_{a \in A} Q^{*}(s,a).
\end{align}
Thus, the agent only needs to know the state-action value function $Q^{*}$ to
determine the optimal policy, and does not need to have information about the
transition and reward probabilities. Using the definition of the state-value
function for $Q^{*}$, $V^{*}(s) = Q^{*}(s, \pi^{*}(s))$ gives a system of
equations called \textit{Bellman equations}:
\begin{align}
    V^{*}(s) = \max_{a\in A} \left\{ \mathbb{E}[r(s, a)] + \gamma\sum_{s'\in S}
    \mathbb{P}[s'|s,a] V^{*}(s') \right\} \quad \forall s \in S.
\end{align}
\begin{theorem}[Bellman Equations]
    The policy value $V_\pi(s)$ of a policy $\pi$ at state $s$ for an MDP
    with infinite horizon is subjected to the following linear system of
    equations
    \begin{align}
        V^{*}(s) = \mathbb{E}_{a_1 \sim \pi(s)}[r(s, a_1)] + \gamma\sum_{s'\in S}
        \mathbb{P}[s'|s, \pi(s)] V^{*}(s') \quad \forall s \in S.
    \end{align}
\end{theorem}
We can rewrite the linear system to
\begin{align}
    v = \gamma P v + r,
\end{align}
where
\begin{align}
    &p_{s, s'} = \mathbb{P}[s'|s, \pi(s)] \quad \forall s,s' \in S \quad
    \text{is the transition probability matrix,}\\
    &v_s = V_\pi(s) \quad \forall s \in S \quad \text{is the value vector and }\\
    &r_s = \mathbb{E}[r(s, \pi(s))] \quad \forall s \in S \quad \text{is the reward
    vector}.
\end{align}
For an MDP with finite horizon with a known probability matrix $P$ and a
known reward expectation $r$, linear algebra gives a unique solution for
the value vector by matrix inversion
\begin{align}
    v^{*} = (I - \gamma P)^{-1} R,
\end{align}
where $I$ is the appropriate identity matrix.
\newline

Proofs for the theorems can be found in \cite[Chapter~17]{ml_foundations}.
\section{Monte Carlo Tree Search (MCTS)}
MCTS is a search algorithm for decision processes like MDP. The algorithm
focuses on finding the most favorable course of action by uniformly expanding
the search tree at each step and applying random actions/decisions until
reaching a terminal state, this outcome is then communicated back to the
start via backpropagation. In the last decade MCTS has been combined with
learning based methods like AlphaZero \cite{alphazero} to produce agents
which can beat humans in Chess, GO and other games.
\newline

On a basic level the concept of the algorithm based on approximating the true
value of an action through random simulation, which ultimately
adjusts the policy efficiently to derive a best-first strategy
\cite{mcts_survey}. One search iteration is divided into four steps:
\newline

\textbf{1. Selection Phase:} Also called the Tree-Traversal Phase. Here we
start at the root node state and walk down until reaching a leaf node state
with a selection policy. The leaf node state is a state which is not yet
fully expanded, meaning it has children that are not visited or it is not a
terminal state. The selection policy of the walk is based on the upper
confidence bound ($UCB1$). The idea is to weight out and balance the
`exploration vs exploitation" trade-off. We select an action, leading to then
next node state, which has the biggest $UCB1$ value. The standard $UCB1$
\cite{sutton_barto} formula is
\begin{align}
    UCB1 = \frac{w_i}{n_i} + C \sqrt{\frac{\ln N_i}{n_i}},
\end{align}
where $w_i$ is the number of wins of the node, $n_i$ the number of visits of
the node, $N_i$ the number of visits of the parent node and $C\in R$ is a
constant, usually chosen to be $C = \sqrt{2}$
\newline

\textbf{2. Node Expansion Phase:} If the node state is not a terminal one, an
extra node is added into the search tree according to the available actions.
\newline

\textbf{3. Simulation Phase}: Also called the Roll-out state. A
simulation/roll-out evolves around the current node selecting actions
randomly until a terminal node state.
\newline

\textbf{4. Backpropagation Phase:} The result of the simulation is used to
update the information in the nodes traversed, such as the outcome/reward of
the terminal state at the end of the simulation and also the visit count of
the nodes.
\begin{figure}[H]
    \centering
    \begin{tikzpicture}[scale=0.8, transform shape]
        \node[align=center, ellipse, draw, ultra thick] (0)  {Start};

        \node[align=center, draw, ultra thick, below=of 0] (1)
            {$s_0$};

        \node[align=center, draw, ultra thick, below=of 1] (2)
            {are we at\\ a leaf node?};

        \node[align=center, draw, ultra thick, below=of 2] (3)
            {move to child\\ with best UCB};

        \node[align=center, draw, ultra thick, right=1.5cm of 2] (4)
            {was the current state\\ visited before ?};

        \node[align=center, ellipse, draw, ultra thick, right=1.5cm of 4] (5)
            {Roll-out};

        \node[align=center, draw, ultra thick, below=of 4] (6)
            {expand tree for\\each available action};

        \node[align=center, draw, ultra thick, below=of 6] (7)
            {new child state};

        \node[align=center,ellipse, draw, ultra thick, below=of 7] (8)
            {Roll-out};

        \draw[->, ultra thick] (0) to (1);
        \draw[->, ultra thick] (1) to (2);
        \draw[->, ultra thick] (2) to node[midway, right] {NO} (3);
        \draw[->, ultra thick] (2) to node[midway, above] {YES} (4);
        \draw[->, ultra thick] (4) to node[midway, above] {YES} (5);
        \draw[->, ultra thick] (4) to node[midway, right] {NO} (6);
        \draw[->, ultra thick] (6) to (7);
        \draw[->, ultra thick] (7) to (8);

        \draw[->, ultra thick, in=180, out=180] (3) to (2);

    \end{tikzpicture}
    \caption{MCTS algorithmic diagram \label{fig: 3}}
\end{figure}
The general algorithmic structure is described in \cite{mcts_survey}, denoted
in \ref{alg: 1}. For clarity purposes the diagram of the algorithm is in
\ref{fig: 3}.
\begin{algorithm}
\caption{MCTS Algorithm}\label{alg: 1}
    \begin{algorithmic}
        \Require state $s_0$
        \While{within computational budget}
            \State $s_k \gets \text{SelectionPolicy}(s_0)$
            \State  $w \gets \text{RollOut}(s_k)$
            \State $\text{BackPropagate}(s_k, w)$
        \EndWhile
        \Ensure $a\left(\text{BestChild}(s_0)\right)$
    \end{algorithmic}
\end{algorithm}
The search starts with an input state $s_0$, and selects the next set of
states based on the selection policy, which is are the states with the
biggest $UCB1$ value. The state $s_k$ is reached after the selection phase
and the expansion, upon which a simulation is run giving the terminal
result/outcome $w$. The outcome is then communicated back from $s_k$ back to the
root node in the path taken to reach $s_k$ from the root node. At the end the
algorithm returns the  action $a$ leading to the best child of $s_0$.
The number of search iterations is usually indicated by a specific search
time or number of search steps given and is denoted by the while statement.
\newline
The MCTS algorithm is a very diverse algorithm in the sence that it allows
for a wide array of variations and improvements, from changing the selection
policy to droping one of the four basics steps and utilizing
learning parameters.

\section{Graph Neural Networks \label{sec: gnn}}
In this section we go through an overview of the later used the Graph Neural
Networks (GNNs) \cite{gnn_overview}. GNNs are Neural Networks (NNs) which
operate over graph structured data. Let us further denote a graph as a tuple
$(E, V)$, where $E$ are edges and $V$ are the nodes of the graph. Also let
$A$ be the adjacency matrix of the graph, where $a_{ij} = 1$ if $(i, j) \in
E$ and $a_{ij} =0$ otherwise. Finally, we denote $X$ to be the set of node
feature.
\newline

There are multiple types of GNNs some more complex than others
\cite{gnn_overview}, but the core principles are taken from the foundational
convolutional NNs. An image can be seen as a grid graph, where each node
corresponds to a pixel and is connected to its neighboring pixels. Simple
convolution in images aggregates pixel by applying a simple operation
(kernel) on its neighborhood, updating the pixels value. Similarly, in a
graph the convolutional GNN will take the information of the neighboring node
features (and edge features) to update the underlying node by applying some
function/transformation on them \ref{fig: 4}.
\begin{figure}[H]
    \centering
    \begin{tikzpicture}
            \node [style=new style 0] (6) at (3.25, 1) {};
            \node [style=new style 0] (7) at (3, -0.75) {};
            \node [style=new style 0] (8) at (1, -2.5) {};
            \node [style=new style 0] (11) at (-1.5, -1.5) {};
            \node [style=1] (12) at (0.75, 0) {$x_a$};
            \node [style=1] (13) at (-1, 0.25) {$x_c$};
            \node [style=1] (14) at (0.75, 1.25) {$x_b$};
            \node [style=1] (15) at (2, 0) {$x_d$};
            \node [style=1] (16) at (0.75, -1.25) {$x_e$};

            \draw [style=new edge style 0] (6) to (7);
            \draw [style=new edge style 0] (15) to (7);
            \draw [style=new edge style 0] (15) to (6);
            \draw [style=new edge style 0] (16) to (8);
            \draw [style=new edge style 0] (11) to (16);
            \draw [style=new edge style 0] (13) to (11);
            \draw [style=new edge style 1] (12) to (16);
            \draw [style=new edge style 1] (12) to (15);
            \draw [style=new edge style 1] (13) to (12);
            \draw [style=new edge style 1] (14) to (12);



            \node [style=new style 0] (21) at (11.25, 1) {};
            \node [style=new style 0] (22) at (11, -0.75) {};
            \node [style=new style 0] (23) at (9, -2.5) {};
            \node [style=new style 0] (24) at (6.5, -1.5) {};
            \node [style=1] (17) at (8.75, 0) {$x_a'$};
            \node [style=new style 0] (25) at (7, 0.25) {};
            \node [style=new style 0] (26) at (8.75, 1.25) {};
            \node [style=new style 0] (27) at (10, 0) {};
            \node [style=new style 0] (28) at (8.75, -1.25) {};

            \draw [style=new edge style 0] (26) to (17);
            \draw [style=new edge style 0] (17) to (25);
            \draw [style=new edge style 0] (17) to (27);
            \draw [style=new edge style 0] (17) to (28);
            \draw [style=new edge style 0] (28) to (24);
            \draw [style=new edge style 0] (28) to (23);
            \draw [style=new edge style 0] (27) to (21);
            \draw [style=new edge style 0] (27) to (22);
            \draw [style=new edge style 0] (21) to (22);
            \draw [style=new edge style 0] (25) to (24);

            \draw [style=new edge style 1, in=120, out=60, dotted, ->] (12)
                to node[midway, above] {$g(x_a, X_{\mathcal{N}_a}) = x_a'$} (17);
    \end{tikzpicture}
    \caption{Convolutional GNNs: simple example \label{fig: 4}}
\end{figure}

For a graph we can define a node's neighborhood, for a node $i \in V$ a
nodes 1-hop neighborhood is
\begin{align}
    \mathcal{N}_i := \left\{ j: (i,j) \in E\; \wedge \; (j, i) \in E
    \right\}.
\end{align}
The features of the 1-hop neighborhood $\mathcal{N}_i$ of the node are
\begin{align}
    X_{\mathcal{N}_i} = \left\{ \left\{ x_j: j \in \mathcal{N}_i \right\}
    \right\} .
\end{align}
Then the convolutional function $g$ operates on the features of the node
itself and on the features of its 1-hop neighborhood $g\left(x_i,
X_{\mathcal{N}_i}\right)$. Suitable are functions that preserve the invariance
or equivariance of the graph. When applying a permutation matrix $P$ we need
to permute the rows and the columns of the adjacency matrix $A$, which
results to $PAP^{T}$, leading to functions satisfying
\begin{align}
    &f(PX, PAP^{T}) = f(X, A) \quad \text{Invariance}\\
    &f(PX, PAP^{T}) = Pf(X, A) \quad \text{Equivariance}
\end{align}
By applying the local function $g$ on all nodes and their neighborhoods we
get
\begin{align}
    f(X, A) =
    \begin{bmatrix}
        g(x_1, X_{\mathcal{N}_1})\\
        \vdots\\
        g(x_n, X_{\mathcal{N}_n})
    \end{bmatrix},
\end{align}
thus $g$ needs to be permutation invariant and not depend on the order in
$X_{\mathcal{N}_i}$. Then $f$ will be equivariant. Conventionally the local
function $g$ is referred to as `diffusion' , `propagation' or `message
passing' and the global function $f$ as a `GNN Layer'
\cite[Chapter~5]{gnn_overview}. The majority of the GNN Layer types can be
distinguished into three types, \textit{Convolutional}, \textit{Attentional}
and \textit{Message-passing} GNN. They are all different in the sense how $g$
incorporates the neighboring features.
\newline
The convolutional type \cite[Chapter~5]{gnn_overview} as described above for
clarification, the features of the neighbors are aggregated additionally with
fixed weights $c_{ij}$ after applying some function $\psi$ on them
\begin{align}
    x_i' = g\left( x_i, \bigoplus_{j\in \mathcal{N}_i} c_{ij} \psi(x_j) \right).
\end{align}
The attention type \cite{gan_velickovic} we use features of neighbors to
apply implicit weights, via attention. The attention coefficient for each
pair is computed via a function $\alpha_{ij} = a(x_i, x_j)$, usually a
softmax function resulting in
\begin{align}
    x_i' = g\left( x_i, \bigoplus_{j\in \mathcal{N}_i} a(x_i, x_j) \psi(x_j) \right).
\end{align}
The message passing type \cite[Chapter~5]{gnn_overview} computes arbitrary
vectors `messages', these are then sent through the edges back to the
original node which gives
\begin{align}
    x_i' = g\left( x_i, \bigoplus_{j\in \mathcal{N}_i} \psi(x_i, x_j) \right).
\end{align}

In practice the functions $g$ and $\psi$ are learnable and are represented by
some appropriate NN structure and the $\bigoplus$ operator is a nonparametric
operation e.g. sum, mean, maximum or a recurrent neural network
\cite[Chapter~5]{gnn_overview}.
\newline

Dealing with real world networks it is piratical to incorporate edge
attributes in the GNN mechanism. For the attention type this is done with
success in the original paper GAT \cite{gan_velickovic}. Given a set of node
features $(h_i)_{i=1}^{N}\subset  \mathbb{R}^{d}$, the GAT layer produces new
features $(h'_i)_{i=1}^{N} \subset \mathbb{R}^{Kd'}$ with dimension $Kd'$. The
computation of the attention coefficient is the following
\begin{align}
    \alpha_{ij} =
    \text{softmax}_j\left(\text{abs}(e_{ij})\right),
\end{align}
where a attention mechanism $a: \mathbb{R}^{d'}\times\mathbb{R}^{d}\to
\mathbb{R}$ compute the attention coefficients through a learnable
matrix $W\in \mathbb{R}^{d'\times d}$
\begin{align}
    e_{ij}^{(k)} = a(Wh_i, Wh_j).
\end{align}
To stabilize the learning process, a multi-head attention mechanism is
employed with $K$ independent attention heads through
\begin{align}
    h_i' = \|_{k=1}^{K} \sigma\left(
        \sum_{j\in\mathcal{N}_i}\alpha_{ij}^{k} W^{k} h_j
    \right),
\end{align}
where $\sigma$ is an activation function, $\alpha_{ij}^{k}$ is normalized attention
coefficient computed by the $k$-th attention head $a^{k}$ with the
corresponding learnable matrix $W^{k}$ and $\|$ is the concatination
operator. The usual choice for $a$ is sigle layer feed forawrdneural network
with LeakyRelU activation function.
\section{REINFORCE}
In the following section we introduce a policy gradient method called
REINFORCE \cite[Chapter~13]{sutton_barto} on a discrete MDP. For convenience
we set the discount factor $\gamma=1$ throughout the explanation but it can
will be later incorporated into the algorithm. The idea is to
have a parametrization of the policy $\pi_\theta$ through some learnable
parameters $\theta$, which is differentiable wrt. $\theta$ and finite in the
state action space. We would like to find the optimal policy by updating the
weights via some performance measure $J(\theta)$ at the start of an episode
$s_0$ through
\begin{align}
    J(\theta) = V_{\pi_\theta}(s_0),
\end{align}
where $V_{\pi_\theta}$ is the policy value function of the parametrazed
policy $\pi_\theta$. We would like to maximize the preformance measure by
going in the estimated direction of the gradient ascent of $J$ wrt. $\theta$,
such that the updates are of the form
\begin{align}
    \theta_{t+1} = \theta_t + \alpha \mathbb{E}[\nabla J(\theta_t)],
\end{align}
where $\alpha$ is the learning rate and expectation is over some appropriate
distribution. To get an appropriate estimation of the gradient, we take a
look at the gradient of the value function, by equation \ref{eq:
value_expectation} we write the policy value as function of the state value
\begin{align}
    \nabla V_{\pi_\theta}(s)
    &= \nabla \mathbb{E}_{\pi_\theta}
    \left[
        Q_{\pi\theta}(s, a)
    \right] \\
    &= \nabla \left(
        \sum_{a}\pi_\theta(s|a) Q_{\pi_\theta}(s, a)
    \right) \\
    &= \sum_{a} \Big(
        \nabla \pi_\theta(a|s)Q_{\pi_\theta}(s,a) + \pi_\theta(a|s)\nabla
        Q_{\pi_\theta}(s, a)
    \Big),
\end{align}
where the gradient here is always wrt. $\theta$. By the definition of the
state-action function \ref{def: state-action} we have
\begin{align}
    \nabla V_{\pi_\theta}(s)
    &= \sum_{a} \Big(
        \nabla \pi_\theta(a|s)Q_{\pi_\theta}(s,a) + \pi_\theta(a|s)\nabla
        \sum_{s',r'} \mathbb{P}[s', r'|s, a]\left( r+V_{\pi_\theta}(s') \right)
    \Big)\\
    &= \sum_{a} \Big(
        \nabla \pi_\theta(a|s)Q_{\pi_\theta}(s,a) + \pi_\theta(a|s)
        \sum_{s'} \mathbb{P}[s'|s, a]
        \nabla V_{\pi_\theta}(s)
    \Big).
\end{align}
We can keep applying the same calculation to the term in $\nabla
V_{\pi_\theta}(s)$ to get an expectation under the probabilities of landing
in some state $x$ from the state $s$ in $k$ steps under $\pi$, i.e.
\begin{align}
    \nabla V_{\pi_\theta}(s)
    &= \sum_{x\in S}\sum_{k \in \mathbb{N} }\mathbb{P}[s\rightarrow x, k,
    \pi]   \sum_{a} \nabla \pi_\theta(a|s)Q_{\pi}(s,a).
\end{align}
Then we have that
\begin{align}
    \nabla J(\theta)
    &= \nabla V_{\pi_\theta}(s_0)\\
    &= \sum_{s}\sum_{k \in \mathbb{N} }\mathbb{P}[s_0\rightarrow s, k,
    \pi]   \sum_{a} \nabla \pi_\theta(a|s)Q_{\pi}(s,a)\\
    &= \sum_{s}\eta(s)  \sum_{a} \nabla \pi_\theta(a|s)Q_{\pi}(s,a)\\
    &= \sum_{s'}\eta(s') \sum_{s} \frac{\eta(s)}{\sum_{s'}\eta(s') }
    \sum_{a} \nabla \pi_\theta(a|s)Q_{\pi}(s,a)\\
    &= \sum_{s'}\eta(s')\sum_{s}\mu(s)\sum_{a} \nabla
    \pi_\theta(a|s)Q_{\pi}(s,a)\\
    &\propto \sum_{s}\mu(s)\sum_{a} \nabla
    \pi_\theta(a|s)Q_{\pi}(s,a).
\end{align}
Where the constant to which the last equation is proportional to is the
average length of an episode. The distribution $\mu$ is the on-policy
distribution over the regular $\pi$ \cite[Chapter~9,10]{sutton_barto}, thus
\begin{align}
    \nabla J(\theta) \propto \mathbb{E}_\pi\left[
        \sum_{a}Q_{\pi_{\theta}}(s_t, a)\nabla\pi_\theta(a|s_t)
    \right]
\end{align}
by replacing $s$ by some sample $s_t$. Doing the same for $a$ to $a_t \sim
\pi$ we can expand by $\pi_\theta$ to preserve equality
\begin{align}
    \nabla J(\theta)
    &= \mathbb{E}_\pi\left[
        \sum_{a}Q_{\pi_{\theta}}(s_t,
        a) \pi_\theta(a|s_t)\frac{\nabla\pi_\theta(a|s_t)}{\pi_\theta (a| s_t)}
    \right]\\
    &= \mathbb{E}_\pi\left[
        Q_{\pi_{\theta}}(s_t,
        a_t) \frac{\nabla\pi_\theta(a_t|s_t)}{\pi_\theta(a_t|s_t)}
    \right]\\
    &= \mathbb{E}_\pi\left[
        G_t \frac{\nabla\pi_\theta(a_t|s_t)}{\pi_\theta(a_t|s_t)}
    \right],
\end{align}
where $G_t$ is the sum of the discounted rewards from $t+1$ to the terminal state $T$
final weights update for the REINFORCE method is
\begin{align}
    \theta_{t+1} = \theta_t + \alpha G_t
    \frac{\nabla \pi_\theta(a_t|s_t)}{\pi_\theta(a_t|s_t)},
\end{align}
where for the expression $\frac{\nabla
\pi_\theta(a_t|s_t)}{\pi_\theta(a_t|s_t)}$, we can just compute the gradient
of natural logarithm, i.e. $\nabla \ln \pi_\theta(a_t | s_t)$. Hence the
reinforce algorithm has the following structure
\begin{algorithm}
\caption{REINFORCE}\label{alg: reinforce}
    \begin{algorithmic}
        \Require Policy parametrization $\pi_\theta$, step-size $\alpha$,
        initial parameters $\theta$
        \While{within computational budget}
            \State Generate an episode $(s_0, a_0), \ldots, (s_T, a_T)$
            through $\pi_\theta$
            \For{$t=0,1,\ldots,T-1$}
                \State $G_t = \sum_{k=t+1}^T\gamma^{k-t-1}  r_k$
                \State $\theta \gets \theta + \alpha\gamma^{t} G_t
                \nabla\ln\pi_\theta(a_t|s_t)$
            \EndFor
        \EndWhile
    \end{algorithmic}
\end{algorithm}
Furthermore we can generalize the gradient updates to include a baseline
function $b(s)$ \cite[Chapter~13.4]{sutton_barto}, instead of updating $G_t$
we centralize the rewards by
\begin{align}
    G_t - b(s_t),
\end{align}
giving a much better performance. This is called the REINFORCE with baseline
algorithm and there are a verity of options when construing the baseline
function, some work better then others.
\section{REINFORCE with MCTS\label{sec: gen}}
In this section we give a brief overview of how we can utilize MCTS with
REINFORCE. Instead of doing a one-per-one episode, we are going
to compute batches of episodes and average out the Monte
Carlo improved policy gradients to compute the updates. There are two
separate phases for making the algorithm work with purely reinforcement
learning, in the one phase the model plays with itself and gathers
information about the environment and gathers the `data'. In the second phase
the information of the self play part is used to update the weights of the
policy parametrization \ref{fig: 5}.
\begin{figure}[H]
    \centering
    \begin{tikzpicture}
        \node[font=\bfseries] (0) {Self-Play};
        \node[right=of 0] (1) {$\pi_{\theta}(s)$};
        \node[font=\bfseries, right= of 1] (2) {Train};

        \draw[->, ultra thick] (0) to[out=90, in=90] node[midway, above]
            {gather episodes} (2);
        \draw[->, ultra thick] (2) to[out=270, in=270] node[midway, below]
            {optimize} (0);
    \end{tikzpicture}
    \caption{Self-Play loop \label{fig: 5}}
\end{figure}
The policy is parametrized through through the parameters $\theta$ of a deep
neural network $\pi_\theta$, which takes a state as input and outputs the
transition probability vector $\pi_\theta(s)$ for each action. In the
self-play phase the model uses the MCTS algorithm to traverse down the MDP
tree from the root node $s_0$. The key difference from traditional MCTS is
the selection strategy, the $UCB$ formula used is
\begin{align}
    UCB = \max_a \left(Q(s,a) + U(s,a)\right), \label{eq: mod_ucb}
\end{align}
where $Q(s,a)$ is the state-action value, and
\begin{align}
U(s,a)=C \pi_\theta(a|s)\sqrt{\frac{\sum_{i}N(s, a_i)}{1+N(s,a)}},
\end{align}
where $C$ is the exploration-exploration constant, $\pi_\theta(a|s)$ is the
probability output for action $a$ at state $s$ of the parametrized policy,
$N(s, a)$ the node visit count and $\sum_{i} N(s, a_i)$ is the visit count of
the parent. Generally it is non-trivial to find a suitable value for the
constant $C$, theoretically it is set to $\sqrt{2}$ in the multi-armed bandit
attack and MDP environments which have rewards spread evenly in the interval
$[0, 1]$. To avoid time-consuming tuning of the constant $C$ we normalize the
$Q$ values during backpropagation. For the normalization we keep track of two
quantities, the $Q$ value of the state action and the best value $b$. The $Q$
value does not exceed $1$ and at the start of each search the best value of
the root node is set to $b_0 = 1$ and else $0$. Continuously during the
backpropagation phase if the best of the child value exceeds the best value
of the parent, then $b_p$ (parent best value) is updated to the better value
and the $Q$ value is then set to
\begin{align}
    Q_c = \frac{r_c}{b_p},
\end{align}
where the subscript $c$ denotes the child and the subscript $p$ the parent
node, $r_c$ is the estinaded reward of the child from the simulation phase
and the current state and $b_p$ is the best value of the parent.
\newline

After an MCTS search we get the so called MCTS improved policy at state $s$
for action $a$,
\begin{align}
    \pi^{\text{MCTS}}(a|s) = \frac{Q(s, a)}{\sum_{a'}Q(s, a')}
\end{align}
We choose the action with the highes probability
\begin{align}
    a = \argmax_a \pi^{\text{MCTS}}(a|s),
\end{align}
transitioning to the next state and doing the MCTS search on the new state.
After playing $N_{\text{play}}$ times, for each time step of each episode we
compute the average of $\nabla \ln \pi_\theta$ and $G_t$ to update the
parameters as in REINFORCE \ref{alg: reinforce}.
